% open-logic.tex
% compiles to produce complete text using open-logic-dev style

\documentclass[../include/open-logic-part]{subfiles}

\begin{document}

\clearpage

\begin{editorial}
Dit bestand haalt alle inhoud van het Open Logic Project binnen. Zie je een
opmerking van de redactie zoals deze, dan is het bestand gemaakt zonder te
letten op hoe het materiaal aan de lezer wordt aangeboden. Als je dit kunt
lezen, kun je deze pdf waarschijnlijk beter \emph{niet} gebruiken om les te
geven of te studeren.

Met allerlei hulpmiddelen van het Open Logic Project kun je een tekst maken die
beter past bij lesgeven of zelfstudie. In de standaardversie gelden alle logische
operatoren bijvoorbeeld als primitief: we omschrijven ze niet met andere
operatoren. In bewijzen werken we dan voor elke operator alle gevallen uit.
Maar sommige van die gevallen kun je veel beter als opgave laten staan. Het Open
Logic Project is ook nog niet af. Om anderen te laten meewerken en de tekst te
verbeteren, staan er ook stukken in die nog niet af zijn of waarin bijvoorbeeld
nog opgaven ontbreken. Een pdf voor een cursus laat zulke
stukken weg.

Voor pdf's die beter werken bij lesgeven en studeren kun je kijken naar de
\href{http://builds.openlogicproject.org/}{voorbeeldcursussen op de
  OLP-website}. Wil je zelf een pdf maken, dan kun je beginnen met het
\href{https://github.com/OpenLogicProject/OpenLogic/tree/master/courses/sample}{voorbeeldbestand
  waarmee je de tekst samenstelt}. In de bronbestanden van de leerboeken die van
dit project zijn afgeleid, vind je uitgebreidere en moeilijkere voorbeelden.
\end{editorial}

\olimport[sets-functions-relations]{sets-functions-relations-complete}

\olimport[propositional-logic]{propositional-logic}

\olimport[first-order-logic]{first-order-logic}

\olimport[model-theory]{model-theory}

\olimport[computability]{computability}

\olimport[turing-machines]{turing-machines}

\olimport[incompleteness]{incompleteness}

\olimport[second-order-logic]{second-order-logic}

\olimport[lambda-calculus]{lambda-calculus}

\olimport[many-valued-logic]{many-valued-logic}

\olimport[normal-modal-logic]{normal-modal-logic}

\olimport[applied-modal-logic]{applied-modal-logic}

\olimport[intuitionistic-logic]{intuitionistic-logic}

\olimport[counterfactuals]{counterfactuals}

\olimport[set-theory]{set-theory}

\olimport[methods]{methods}

\olimport[history]{history}

\olimport[reference]{reference}
\end{document}

